\documentclass[11pt]{article}
\usepackage{url}
\usepackage{alltt}
\usepackage{bm}
\linespread{1}
\textwidth 6.5in
\oddsidemargin 0.in
\addtolength{\topmargin}{-1in}
\addtolength{\textheight}{2in}

\usepackage{amsmath}
\usepackage{amssymb}
\usepackage{hyperref}

\begin{document}


\begin{center}
\Large
STA 214 Homework 2\\
\normalsize
\vspace{5mm}
\end{center}

\noindent \textbf{Due:} Friday, September 11, 11:59pm on Canvas.\\ 

\noindent \textbf{Instructions:} There are two parts to this assignment. Part I is practice with logistic regression,  and Part II is practice with maximum likelihood estimation.\\

\noindent \textbf{Getting started:} Begin by downloading the HW2 template from the course website:\\

\url{https://sta214-f26.github.io/homework/hw_02_template.qmd}\\

\noindent Save this template file to your computer, then open it in RStudio. As you complete the assignment, you will write down your answers to all questions in the file, and include all R code in code chunks. \textit{If a question requires code, you will not receive credit if no code is provided.}\\

\noindent \textbf{Collaborators and AI statement:} At the top of your assignment, list any collaborators (Dr. Evans does not count -- he is a resource, not a collaborator), and provide an AI statement describing your use of AI to complete the questions. If you did not use AI in any way, let me know that too. If you did use AI, provide: the model used, the HW questions on which you used AI, and a summary of the ways in which AI provided assistance. As described in the syllabus, AI use is permitted on homework assignments, but submissions without a complete AI statement may not receive full credit.\\

\noindent \textbf{Submission:} When you have completed the assignment, render your homework to HTML and submit on Canvas. For Part II, you may submit a separate scan of your written work, if you prefer.

\section*{Data}

The RMS Titanic was a huge, luxury passenger liner designed and built in the early 20th century. Despite the fact that the ship was believed to be unsinkable, during her maiden voyage on April 15, 1912, the Titanic collided with an iceberg and sank. Of all the passengers and crew, less than half survived. Part of the reason why so few people survived has been attributed to the fact that the Titanic did not carry enough lifeboats for its passengers and crew. This meant that there was competition for space in the boats, and not everyone was able to make it aboard. Communication errors, stress and shock...there were a great many factors that contributed to this tragedy.\\

\noindent The loss of life during the Titanic tragedy was enormous, but there were survivors. Was it random chance that these particular people survived? Or were there some specific characteristics of these people that led to their positions in the life boats? Let's investigate.\\

\noindent We have observations on 12 different variables, some categorical and some numeric:

\begin{itemize}
\item \verb;Passenger;: A unique ID number for each passenger.
\item \verb;Survived;: An indicator for whether the passenger survived (1) or perished (0) during the disaster.
\item \verb;Pclass;: Indicator for the class of the ticket held by this passengers; 1 = 1st class, 2 = 2nd class, 3 = 3rd class.
\item \verb;Name;: The name of the passenger.
\item \verb;Sex;: Binary indicator for the sex of the passenger.
\item \verb;Age;: Age of the passenger in years; Age is fractional if the passenger was less than 1 year old.
\item \verb;SibSp;: number of siblings/spouses the passenger had aboard the Titanic. Here, siblings are defined as brother, sister, stepbrother, and stepsister. Spouses are defined as husband and wife.
\item \verb;Parch;: number of parents/children the passenger had aboard the Titanic. Here, parent is defined as mother/father and child is defined as daughter,son, stepdaughter or stepson. NOTE: Some children traveled only with a nanny, therefore parch=0 for them. There were no parents aboard for these children.
\item \verb;Ticket;: The unique ticket number for each passenger.
\item \verb;Fare;: How much the ticket cost in US dollars.
\item \verb;Cabin;: The cabin number assigned to each passenger. Some cabins hold more than one passenger.
\item \verb;Embarked;: Port where the passenger boarded the ship; C = Cherbourg, Q = Queenstown, S = Southampton
\end{itemize}

\noindent \textbf{Goal:} Our goal is to predict the probability that a passenger survives the Titanic disaster.

\subsection*{Loading the data}

The \verb;titanic; data can be loaded into R with the following command:

\begin{verbatim}
titanic <- read.csv("https://sta214-f26.github.io/homework/Titanic.csv")
\end{verbatim}

\noindent Here \texttt{read.csv} is a function that imports data from a CSV file. We can pass \texttt{read.csv} either a local path on our computer, or a URL -- in this case, we use the URL where the data is stored online. We have called the data \verb;titanic; in R.\\

\noindent Copy the command to load the data into the setup chunk of your Quarto file, and run it.

\newpage

\section{Logistic regression}

In the first part of the assignment, we will fit a logistic regression model to the Titanic data. If you get stuck, remember that the course codebook may be helpful: \url{https://sta214-f26.github.io/resources/codebook.html} 

\begin{enumerate}
\item Before we fit any models, we need to do a bit of data exploration. To keep this assignment manageable, we will only do the first few steps here, but a full analysis would require additional EDA.

\begin{enumerate}
\item Looking at the available variables, there are some that are not valid choices for explanatory variables, meaning we can not use them as X variables in a parametric model. Which are they, and why can they not be used?

\item Does the data set have any missing data? If so, in which variables?

\item Create a table showing how many passengers in the data survived and perished.
\end{enumerate}

\item Next, we want to fit a logistic regression model to predict the probability that a patient survives the disaster. For this question, use passenger class, sex, and age as the explanatory variables.

\begin{enumerate}
\item Write down the population logistic regression model, making sure to include both the random and systematic components. Use proper notation and include all subscripts. (See ``writing math in Quarto and R Markdown'' in the course codebook, \url{https://sta214-f26.github.io/resources/codebook.html}) \textbf{Note:} passenger class is a categorical variable with more than two levels, but is recorded as a number in the data. Use \texttt{as.factor()} to convert passenger class into a categorical variable.

\item In R, fit your logistic regression model, and write down the equation of the fitted model. (See ``writing math in Quarto and R Markdown'' in the course codebook, \url{https://sta214-f26.github.io/resources/codebook.html})

\item Interpret each estimated coefficient (in terms of the odds).
\end{enumerate}

\item Finally, let's use the logistic regression model to make some predictions!

\begin{enumerate}
\item Suppose we have a 20 year old male passenger, but we don't know their passenger class. Which passenger class (first, second, or third) would have the highest chance of survival? You should answer the question without doing any calculations, and explain your reasoning.

\item Suppose we have a 30 year old passenger, but we don't know their passenger class or sex. Which combination of class and sex would have the highest chance of survival? You should answer the question without doing any calculations, and explain your reasoning.

\item \textit{In the observed data}, what is the highest predicted probability of survival for a female first-class passenger? (Hint: the \verb;range; function in R will give you the highest and lowest values of a variable or vector)

\end{enumerate}

\end{enumerate}

\newpage

\section{Maximum likelihood estimation}

This part focuses on practice with maximum likelihood estimation, and is separate from the Titanic data you were working with above in Part I. A brief review of mathematical notation and rules for logs and derivatives can be found in the course codebook.\\

\noindent Suppose we have a random variable $Y$, which can take on values $y = 0, 1, 2, 3, 4,...$ (i.e., any non-negative integer). The probability of each outcome does not depend on any explanatory variables, and we are told that
$$P(Y = y | \lambda) = \dfrac{\lambda^y e^{-\lambda}}{y!}$$
where $\lambda > 0$ is an unknown parameter. We want to estimate $\lambda$.\\

\noindent We get $n$ independent observations $Y_1,...,Y_n$ from this distribution, and we will use the method of maximum likelihood to estimate $\lambda$.

\begin{enumerate}
\item[4.] 

\begin{enumerate}
\item Write down the likelihood $L(\lambda)$.

\item Using calculus, show that the maximum likelihood estimate is $\widehat{\lambda} = \frac{1}{n} \sum \limits_{i=1}^n Y_i$. Make sure to show all your work.
\end{enumerate}
\end{enumerate}

\end{document}
